%=====================================================================
% Datos
%   Una subsección por archivo en datos/. Las tablas de registros las
%   genera bd/exportar_datos.py ejecutando bd/consultas_datos.sql sobre la
%   base cargada; no se editan a mano.
%=====================================================================
\section*{Datos}
\addcontentsline{toc}{section}{Datos}

% Holgura para las líneas que terminan en una ruta de archivo o un código largo.
\setlength{\emergencystretch}{2em}

% \captura{archivo}{pie}: una captura de datos/capturas, a lo ancho, con su pie.
\newcommand{\captura}[2]{%
  \par\medskip\noindent\begin{minipage}{\linewidth}\centering
  \includegraphics[width=\linewidth]{datos/capturas/#1}\par\smallskip
  {\small\textit{#2}}
  \end{minipage}\par\medskip}

% \consulta{nombre}: el texto de una consulta de bd/consultas_datos.sql.
\newcommand{\consulta}[1]{\lstinputlisting[style=sql, basicstyle=\ttfamily\scriptsize]{bd/tex/datos/#1.sql}}

%=====================================================================
% Datos · Registros cargados
%=====================================================================
\subsection*{Registros cargados}
\addcontentsline{toc}{subsection}{Registros cargados}

Los registros con que se probó la base de datos son los que carga \at{bd/insertar\_datos\_prueba.sql}, el script de la sección anterior: \dpNumFilas{} registros repartidos en las \dpNumTablasConDatos{} tablas del modelo, muchos más que los veinte que se piden como mínimo. No son filas de relleno. Forman un escenario del juego, fotografiado el \dpFecha{} a las 18:00 UTC, en el que cada registro representa una situación que describen los casos de uso: cuentas en cada estado, partidas de los tres tipos terminadas de las cinco formas, compras, cajas, reportes y sanciones. Cada cosa que los datos deben permitir probar tiene su subsección:

\begin{center}\small
\begin{longtable}{|L{0.22\textwidth}|L{0.2\textwidth}|L{0.48\textwidth}|}
\hline
\textbf{Qué se prueba} & \textbf{Dónde} & \textbf{Con qué registros} \\ \hline
\endfirsthead
\hline
\textbf{Qué se prueba} & \textbf{Dónde} & \textbf{Con qué registros} \\ \hline
\endhead
Casos normales & \textit{Casos normales} & Una partida clasificatoria de principio a fin y las compras en la tienda, con todo lo que dejan escrito en otras tablas. \\ \hline
Relaciones entre tablas & \textit{Relaciones entre tablas} & Consultas que recorren las llaves foráneas: de la partida a sus jugadores, de la cuenta a lo que lleva equipado, del reporte a la sanción y a la apelación, de la caja a su contenido, y de una sanción a la que la sustituye. \\ \hline
Diferentes estados o categorías & \textit{Estados y categorías} & \dnNumValoresConFilas{} de los \dnNumValores{} valores posibles de las \dnNumDominios{} columnas de estado o de categoría tienen al menos un registro; los que faltan se explican uno por uno. \\ \hline
Variaciones de la información & \textit{Variaciones} & Nulos que dependen del tipo de fila, cuentas eliminadas y baneadas, textos en inglés y con emojis, partidas sin reloj, de cuatro jugadores y perdidas por tiempo, jugadas deshechas, objetos retirados. \\ \hline
Que la base cumple sus reglas & \textit{Casos que la base rechaza} & \dnNumPruebasRechazo{} operaciones que el documento prohíbe, intentadas sobre estos registros; la base rechaza las \dnNumPruebasRechazadas{}. \\ \hline
Todos los registros & \textit{Registros por tabla} & El contenido de cada tabla, al final de la sección. \\ \hline
\end{longtable}
\end{center}

\textbf{De dónde salen las tablas.} Ninguna tabla de registros de esta sección se copió a mano: son el resultado de las consultas de \at{bd/consultas\_datos.sql} sobre la base cargada, que \at{bd/exportar\_datos.py} ejecuta y convierte en tablas del documento. Las mismas consultas pueden abrirse en SQL Server Management Studio y ejecutarse una a una. Para regenerar las tablas después de cambiar el escenario:

\begin{lstlisting}[style=sql]
python bd\generar_datos_prueba.py
sqlcmd -S localhost -E -b -I -f 65001 -i bd\insertar_datos_prueba.sql
python bd\exportar_datos.py
\end{lstlisting}

\textbf{Las capturas.} Se tomaron el 13 de septiembre de 2026 en la consola de Windows, ejecutando \at{sqlcmd} sobre SQL Server 2025 (17.0.4065.4). Cada una muestra el comando y su salida tal como aparecieron; la consola no tiene fuente para el emoji de un mensaje de chat y lo dibuja como un recuadro, aunque la base lo guarda completo. La captura 1 es la carga: los tres scripts, en orden, sobre la instancia, y las \dpNumComprobaciones{} comprobaciones del final de \at{insertar\_datos\_prueba.sql}, todas con cero incoherencias.

\captura{01-carga}{Captura 1. Creación de la base y de las tablas, carga de los registros y comprobaciones de coherencia.}

\textbf{Registros por tabla.} La tabla siguiente, que genera el script de datos, cuenta los registros de cada tabla por área; la captura 2 los cuenta en la propia base, consultando el catálogo del sistema, y da las mismas cifras.

% Archivo generado por bd/generar_datos_prueba.py. No editar a mano.
\begin{center}\small
\begin{longtable}{|L{0.34\textwidth}|L{0.12\textwidth}|}
\hline
\textbf{Tabla} & \textbf{Filas} \\ \hline
\endfirsthead
\hline
\textbf{Tabla} & \textbf{Filas} \\ \hline
\endhead
\multicolumn{2}{|l|}{\textbf{Catálogos precargados (\mbox{D-09})}} \\ \hline
\textit{Ranura} & 6 \\ \hline
\textit{ObjetoCosmetico} & 23 \\ \hline
\textit{Modo} & 3 \\ \hline
\textit{Division} & 4 \\ \hline
\textit{NivelIA} & 4 \\ \hline
\textit{Leccion} & 7 \\ \hline
\textit{TipoCaja} & 3 \\ \hline
\textit{TipoCajaObjeto} & 14 \\ \hline
\textit{PalabraProhibida} & 8 \\ \hline
\textit{MotivoReporte} & 5 \\ \hline
\multicolumn{2}{|l|}{\textbf{Cuenta y acceso}} \\ \hline
\textit{Usuario} & 11 \\ \hline
\textit{Sesion} & 15 \\ \hline
\textit{CodigoSegundoFactor} & 4 \\ \hline
\textit{TokenVerificacionCorreo} & 12 \\ \hline
\textit{TokenRecuperacion} & 3 \\ \hline
\textit{AceptacionTerminos} & 11 \\ \hline
\multicolumn{2}{|l|}{\textbf{Perfil y personalización}} \\ \hline
\textit{HistorialNickname} & 1 \\ \hline
\textit{Avatar} & 3 \\ \hline
\textit{EnlaceRed} & 2 \\ \hline
\textit{UsuarioObjeto} & 95 \\ \hline
\textit{Equipamiento} & 66 \\ \hline
\multicolumn{2}{|l|}{\textbf{Partida}} \\ \hline
\textit{Partida} & 12 \\ \hline
\textit{Participacion} & 25 \\ \hline
\textit{ParticipacionObjeto} & 150 \\ \hline
\textit{Jugada} & 134 \\ \hline
\textit{OfertaTablas} & 4 \\ \hline
\textit{EnlaceEspectador} & 1 \\ \hline
\multicolumn{2}{|l|}{\textbf{Emparejamiento, salas y chat}} \\ \hline
\textit{ColaEmparejamiento} & 1 \\ \hline
\textit{Sala} & 2 \\ \hline
\textit{SalaParticipante} & 2 \\ \hline
\textit{Invitacion} & 4 \\ \hline
\textit{Mensaje} & 11 \\ \hline
\multicolumn{2}{|l|}{\textbf{Clasificación y economía}} \\ \hline
\textit{EstadisticaModo} & 33 \\ \hline
\textit{HistorialDivision} & 2 \\ \hline
\textit{Caja} & 3 \\ \hline
\textit{CajaContenido} & 3 \\ \hline
\textit{MovimientoMoneda} & 30 \\ \hline
\multicolumn{2}{|l|}{\textbf{Relaciones entre jugadores y tutorial}} \\ \hline
\textit{Amistad} & 4 \\ \hline
\textit{Solicitud} & 1 \\ \hline
\textit{Silencio} & 1 \\ \hline
\textit{Bloqueo} & 1 \\ \hline
\textit{ProgresoTutorial} & 8 \\ \hline
\multicolumn{2}{|l|}{\textbf{Moderación y bitácoras}} \\ \hline
\textit{Reporte} & 5 \\ \hline
\textit{Sancion} & 5 \\ \hline
\textit{Apelacion} & 3 \\ \hline
\textit{BitacoraModeracion} & 19 \\ \hline
\textit{BitacoraAcceso} & 54 \\ \hline
\textbf{Total} & \textbf{818} \\ \hline
\end{longtable}
\end{center}


\captura{02-conteo}{Captura 2. Registros de cada tabla según SQL Server.}

%=====================================================================
% Datos · Escenario
%   Qué representa cada cuenta y cada partida de los datos de prueba.
%=====================================================================
\subsection*{Escenario}
\addcontentsline{toc}{subsection}{Escenario}

Los registros cuentan una historia coherente: cada cuenta y cada partida está ahí para representar una situación de los casos de uso, y todo lo que el modelo guarda como consecuencia ---el elo, las estadísticas, el saldo, el estado del tablero--- se calculó a partir de ellas, como se explica en la sección de implementación.

\subsubsection*{Cuentas}
\addcontentsline{toc}{subsubsection}{Cuentas}

Las \dpNumUsuarios{} cuentas cubren todos los tipos, estados y roles de \textit{Usuario}:

\begin{center}\small
\begin{longtable}{|L{0.04\textwidth}|L{0.18\textwidth}|L{0.2\textwidth}|L{0.48\textwidth}|}
\hline
\textbf{Id} & \textbf{Nickname} & \textbf{Tipo, estado y rol} & \textbf{Qué representa} \\ \hline
\endfirsthead
\hline
\textbf{Id} & \textbf{Nickname} & \textbf{Tipo, estado y rol} & \textbf{Qué representa} \\ \hline
\endhead
1 & \at{admin\_bastion} & Registrada, activa, administrador & Otorgó el rol de moderadora y resolvió dos apelaciones, una aceptada y otra rechazada. Tres contraseñas incorrectas le bloquearon el acceso cinco minutos (\mbox{CU-01} \mbox{RN-03}); después pidió dos veces el código de recuperación y el segundo invalidó al primero (\mbox{CU-03} \mbox{RN-05}). \\ \hline
2 & \at{ana\_modera} & Registrada, activa, moderadora & Resolvió reportes con sanción y sin ella, baneó una cuenta y sustituyó una sanción por otra. Dejó que un reporte volviera a la cola y tiene otro en revisión. \\ \hline
3 & \at{luis\_quo} & Registrada, activa, jugador & Seis partidas clasificatorias en el modo clásico, división plata y racha de tres; compró objetos y una caja; tiene un cambio de correo pendiente y juega la partida en curso. \\ \hline
4 & \at{marta\_muros} & Registrada, activa, jugador & Segundo factor activado, con un código pendiente porque entra desde otro dispositivo (\mbox{D-01}); ganó por tiempo una partida rápida; espera en la cola; una caja de recompensa caducó sin abrirse. \\ \hline
5 & \at{pedro\_peon} & Registrada, activa, jugador & Partida contra la IA con un deshacer; perdió por tiempo una partida rápida; anfitrión de la sala abierta; una sanción de chat vencida, cuya apelación se rechazó. \\ \hline
6 & \at{sofia\_salto} & Registrada, activa, jugador & Usa la interfaz en inglés, aceptó los términos en ese idioma y escribe en inglés en el chat. Cambió su nickname una vez, recuperó su contraseña y una apelación aceptada le retiró una sanción. \\ \hline
7 & \at{Invitado\_4821} & Invitado, activa, jugador & Sin correo ni contraseña; aceptó una invitación, jugó una partida privada y está en la sala abierta. \\ \hline
8 & \at{nico\_nuevo} & Registrada, pendiente, jugador & No verificó su correo: tiene un token vigente y otro invalidado por el reenvío. \\ \hline
9 & \at{rayo\_veloz} & Registrada, suspendida, jugador & Reportado y suspendido; su suspensión se amplió con una segunda sanción y apeló. \\ \hline
10 & \at{anonimo\_10} & Registrada, eliminada, jugador & Dada de baja y anonimizada en la misma transacción (\mbox{D-17}). \\ \hline
11 & \at{4dm1n\_oficial} & Registrada, baneada, jugador & Su nickname esquivó el filtro de palabras imitando al administrador; un reporte terminó en una sanción permanente de cuenta (\mbox{CU-44}), y su intento de volver a entrar quedó en la bitácora. \\ \hline
\end{longtable}
\end{center}

\subsubsection*{Partidas}
\addcontentsline{toc}{subsubsection}{Partidas}

Las \dpNumPartidas{} partidas, con \dpNumJugadas{} jugadas, recorren los tres tipos y las cinco formas de término (\mbox{CU-25} \mbox{RN-01}):

\begin{center}\small
\begin{longtable}{|L{0.05\textwidth}|L{0.25\textwidth}|L{0.19\textwidth}|L{0.40\textwidth}|}
\hline
\textbf{Id} & \textbf{Tipo y modo} & \textbf{Jugadores} & \textbf{Cómo termina} \\ \hline
\endfirsthead
\hline
\textbf{Id} & \textbf{Tipo y modo} & \textbf{Jugadores} & \textbf{Cómo termina} \\ \hline
\endhead
1 a 5 & Clasificatorias, clásico & luis y marta & Cuatro a la meta y una rendición. Tras la quinta, los dos entran en una división (\mbox{CU-25} \mbox{RN-14}). \\ \hline
6 & Clasificatoria, rápida & pedro y sofía & Tablas aceptadas, tras una oferta rechazada (\mbox{CU-22}). \\ \hline
7 & Clasificatoria, cuatro jugadores & luis, marta, pedro y sofía & Tres abandonan uno tras otro y gana el que queda (\mbox{CU-25} \mbox{FA-04}). \\ \hline
8 & Contra la IA, clásico, sin reloj & pedro & Se rinde; antes deshizo una jugada, que queda marcada y no se borra (\mbox{D-12}). \\ \hline
9 & Clasificatoria, clásico & luis y rayo & Luis llega a la meta; rayo lo insulta en el chat y es reportado. \\ \hline
10 & Privada, rápida, desde una sala & sofía e invitado & Llegada a la meta. La anfitriona eligió ocho muros (\mbox{CU-18} \mbox{RN-04}) y el invitado entró aceptando su invitación (\mbox{CU-32} \mbox{RN-05}). \\ \hline
11 & Clasificatoria, rápida & marta y pedro & Pedro pierde por tiempo: su reloj llega a cero en su turno, antes de jugar (\mbox{CU-20} \mbox{FA-06}). \\ \hline
12 & Clasificatoria, clásico & luis y sofía & En curso, con una oferta de tablas pendiente y un enlace de espectador. \\ \hline
\end{longtable}
\end{center}

Las jugadas usan la notación del tablero: columnas de la \at{a} a la \at{i} y filas del 1 al 9 (de la \at{a} a la \at{g} y del 1 al 7 en el modo rápido). Un muro se nombra por el surco cuya esquina superior derecha es su centro, más su orientación.

%=====================================================================
% Datos · Casos normales
%=====================================================================
\subsection*{Casos normales}
\addcontentsline{toc}{subsection}{Casos normales}

Un caso normal es el que un caso de uso describe en su flujo normal, sin alternativas ni excepciones. La mayoría de los registros lo son: cuentas registradas y activas que juegan partidas clasificatorias y compran objetos. Dos ejemplos bastan para ver todo lo que un caso normal deja escrito.

\subsubsection*{Una partida clasificatoria completa}
\addcontentsline{toc}{subsubsection}{Una partida clasificatoria completa}

La partida 1 enfrenta a \at{luis\_quo} y \at{marta\_muros} en el modo clásico, con reloj de cinco minutos. Esta consulta reúne la partida, sus dos participaciones y las monedas que cada uno recibió al terminar:

\consulta{normalpartida}
% Archivo generado por bd/exportar_datos.py. No editar a mano.
\begin{center}\scriptsize\setlength{\tabcolsep}{2pt}
\begin{longtable}{|L{46.3pt}|L{57.2pt}|L{42.1pt}|L{62.4pt}|L{59.2pt}|L{41.0pt}|L{29.8pt}|L{38.2pt}|L{55.2pt}|}
\hline
\textbf{\hspace{0pt}Partida} & \textbf{\hspace{0pt}Modo} & \textbf{\hspace{0pt}Forma} & \textbf{\hspace{0pt}Jugador} & \textbf{\hspace{0pt}Resultado} & \textbf{\hspace{0pt}Elo inicial} & \textbf{\hspace{0pt}Elo final} & \textbf{\hspace{0pt}Diferencia} & \textbf{\hspace{0pt}Monedas} \\ \hline
\endfirsthead
\hline
\textbf{\hspace{0pt}Partida} & \textbf{\hspace{0pt}Modo} & \textbf{\hspace{0pt}Forma} & \textbf{\hspace{0pt}Jugador} & \textbf{\hspace{0pt}Resultado} & \textbf{\hspace{0pt}Elo inicial} & \textbf{\hspace{0pt}Elo final} & \textbf{\hspace{0pt}Diferencia} & \textbf{\hspace{0pt}Monedas} \\ \hline
\endhead
1 & CLASICO & META & luis\_\allowbreak{}quo & GANADA & 1000 & 1016 & 16 & 100 \\ \hline
1 & CLASICO & META & marta\_\allowbreak{}muros & PERDIDA & 1000 & 984 & -16 & 25 \\ \hline
\end{longtable}
\end{center}


Los dos empezaron con el elo inicial de 1000 (\mbox{CU-02} \mbox{RN-06}). Con elos iguales, la victoria vale 16 puntos y la derrota los mismos 16 en contra, y la diferencia la calcula la columna \at{diferencia\_elo}. Cada participante recibe un movimiento de monedas de tipo \at{PARTIDA}: 100 el ganador y 25 el perdedor (\mbox{CU-25} \mbox{RN-08}). Las quince jugadas de la partida alternan los turnos, mezclan movimientos y muros y terminan con luis en la fila 9, la meta del jugador que sale de la fila 1:

\consulta{normaljugadas}
% Archivo generado por bd/exportar_datos.py. No editar a mano.
\begin{center}\scriptsize\setlength{\tabcolsep}{3pt}
\begin{longtable}{|L{0.052\linewidth}|L{0.120\linewidth}|L{0.109\linewidth}|L{0.052\linewidth}|L{0.061\linewidth}|L{0.044\linewidth}|L{0.109\linewidth}|L{0.055\linewidth}|L{0.251\linewidth}|}
\hline
\textbf{Número} & \textbf{Jugador} & \textbf{Tipo} & \textbf{Origen} & \textbf{Destino} & \textbf{Surco} & \textbf{Orientación} & \textbf{Tiempo (ms)} & \textbf{Fecha} \\ \hline
\endfirsthead
\hline
\textbf{Número} & \textbf{Jugador} & \textbf{Tipo} & \textbf{Origen} & \textbf{Destino} & \textbf{Surco} & \textbf{Orientación} & \textbf{Tiempo (ms)} & \textbf{Fecha} \\ \hline
\endhead
1 & luis\_\allowbreak{}quo & MOVIMIENTO & e1 & e2 & \textit{NULL} & \textit{NULL} & 7148 & 2026-07-05 17:00:07.148 \\ \hline
2 & marta\_\allowbreak{}muros & MURO & \textit{NULL} & \textit{NULL} & a7 & HORIZONTAL & 6067 & 2026-07-05 17:00:13.215 \\ \hline
3 & luis\_\allowbreak{}quo & MOVIMIENTO & e2 & e3 & \textit{NULL} & \textit{NULL} & 4986 & 2026-07-05 17:00:18.201 \\ \hline
4 & marta\_\allowbreak{}muros & MOVIMIENTO & e9 & d9 & \textit{NULL} & \textit{NULL} & 3905 & 2026-07-05 17:00:22.106 \\ \hline
5 & luis\_\allowbreak{}quo & MOVIMIENTO & e3 & e4 & \textit{NULL} & \textit{NULL} & 2824 & 2026-07-05 17:00:24.930 \\ \hline
6 & marta\_\allowbreak{}muros & MURO & \textit{NULL} & \textit{NULL} & g7 & HORIZONTAL & 10743 & 2026-07-05 17:00:35.673 \\ \hline
7 & luis\_\allowbreak{}quo & MOVIMIENTO & e4 & e5 & \textit{NULL} & \textit{NULL} & 9662 & 2026-07-05 17:00:45.335 \\ \hline
8 & marta\_\allowbreak{}muros & MOVIMIENTO & d9 & d8 & \textit{NULL} & \textit{NULL} & 8581 & 2026-07-05 17:00:53.916 \\ \hline
9 & luis\_\allowbreak{}quo & MOVIMIENTO & e5 & e6 & \textit{NULL} & \textit{NULL} & 7500 & 2026-07-05 17:01:01.416 \\ \hline
10 & marta\_\allowbreak{}muros & MURO & \textit{NULL} & \textit{NULL} & b3 & VERTICAL & 6419 & 2026-07-05 17:01:07.835 \\ \hline
11 & luis\_\allowbreak{}quo & MOVIMIENTO & e6 & e7 & \textit{NULL} & \textit{NULL} & 5338 & 2026-07-05 17:01:13.173 \\ \hline
12 & marta\_\allowbreak{}muros & MOVIMIENTO & d8 & d7 & \textit{NULL} & \textit{NULL} & 4257 & 2026-07-05 17:01:17.430 \\ \hline
13 & luis\_\allowbreak{}quo & MOVIMIENTO & e7 & e8 & \textit{NULL} & \textit{NULL} & 3176 & 2026-07-05 17:01:20.606 \\ \hline
14 & marta\_\allowbreak{}muros & MURO & \textit{NULL} & \textit{NULL} & h2 & HORIZONTAL & 11095 & 2026-07-05 17:01:31.701 \\ \hline
15 & luis\_\allowbreak{}quo & MOVIMIENTO & e8 & e9 & \textit{NULL} & \textit{NULL} & 10014 & 2026-07-05 17:01:41.715 \\ \hline
\end{longtable}
\end{center}


En un movimiento, la jugada tiene casilla de origen y de destino y no tiene surco; en un muro, al revés. Es la regla de subtipo que hace cumplir la restricción \at{CK\_Jugada\_tipo}. La fecha de cada jugada lleva milisegundos, y el tiempo consumido es el que se descuenta del reloj del jugador.

\subsubsection*{Compras en la tienda}
\addcontentsline{toc}{subsubsection}{Compras en la tienda}

Una compra deja dos registros en la misma transacción: un movimiento de monedas de tipo \at{COMPRA}, con importe negativo, y el objeto en el inventario de la cuenta, con origen \at{COMPRA} (\mbox{CU-38} \mbox{RN-08}). La consulta los une y comprueba además si el objeto quedó equipado:

\consulta{normalcompra}
% Archivo generado por bd/exportar_datos.py. No editar a mano.
\begin{center}\scriptsize\setlength{\tabcolsep}{2pt}
\begin{longtable}{|L{44.4pt}|L{54.9pt}|L{92.8pt}|L{43.9pt}|L{43.3pt}|L{74.4pt}|L{50.8pt}|L{31.3pt}|}
\hline
\textbf{\hspace{0pt}Movimiento} & \textbf{\hspace{0pt}Usuario} & \textbf{\hspace{0pt}Objeto} & \textbf{\hspace{0pt}Rareza} & \textbf{\hspace{0pt}Importe} & \textbf{\hspace{0pt}Fecha} & \textbf{\hspace{0pt}Origen en UsuarioObjeto} & \textbf{\hspace{0pt}Equipado} \\ \hline
\endfirsthead
\hline
\textbf{\hspace{0pt}Movimiento} & \textbf{\hspace{0pt}Usuario} & \textbf{\hspace{0pt}Objeto} & \textbf{\hspace{0pt}Rareza} & \textbf{\hspace{0pt}Importe} & \textbf{\hspace{0pt}Fecha} & \textbf{\hspace{0pt}Origen en UsuarioObjeto} & \textbf{\hspace{0pt}Equipado} \\ \hline
\endhead
10 & luis\_\allowbreak{}quo & EMOTE\_\allowbreak{}RISA & RARO & -120 & 2026-07-21 18:00 & COMPRA & sí \\ \hline
13 & luis\_\allowbreak{}quo & PEON\_\allowbreak{}CRISTAL & RARO & -250 & 2026-07-27 18:00 & COMPRA & sí \\ \hline
21 & marta\_\allowbreak{}muros & MURO\_\allowbreak{}LADRILLO & COMUN & -150 & 2026-08-16 19:00 & COMPRA & sí \\ \hline
23 & pedro\_\allowbreak{}peon & EMOTE\_\allowbreak{}RISA & RARO & -120 & 2026-08-21 17:00 & COMPRA & sí \\ \hline
\end{longtable}
\end{center}


El importe de cada compra es el precio del objeto en el catálogo, y no el que envía el cliente (\mbox{CU-38} \mbox{RN-03}). Los cuatro objetos comprados quedaron equipados en su ranura (\mbox{CU-14}).

%=====================================================================
% Datos · Relaciones entre tablas
%=====================================================================
\subsection*{Relaciones entre tablas}
\addcontentsline{toc}{subsection}{Relaciones entre tablas}

Cada consulta de esta subsección une tablas por sus llaves foráneas. Si un registro apuntara a otro que no existe, la unión lo perdería y la consulta devolvería menos filas de las que tiene la tabla; la base, además, no lo permite, como muestra la prueba 5 de los casos que la base rechaza.

\subsubsection*{La partida y sus participantes}
\addcontentsline{toc}{subsubsection}{La partida y sus participantes}

\textit{Partida} se relaciona con \textit{Modo} (N:1) y con \textit{Usuario} a través de \textit{Participacion}, la entidad asociativa que guarda el resultado y el elo de cada jugador. La consulta devuelve las \dpNumPartidas{} partidas con sus participantes: dos en las de dos jugadores, cuatro en la de cuatro y uno en la de la IA, que no es un usuario.

\consulta{relpartidas}
% Archivo generado por bd/exportar_datos.py. No editar a mano.
\begin{center}\scriptsize\setlength{\tabcolsep}{2pt}
\begin{longtable}{|L{35.0pt}|L{82.4pt}|L{77.6pt}|L{60.8pt}|L{44.2pt}|L{64.0pt}|L{44.7pt}|L{27.1pt}|}
\hline
\textbf{\hspace{0pt}Partida} & \textbf{\hspace{0pt}Tipo} & \textbf{\hspace{0pt}Modo} & \textbf{\hspace{0pt}Estado} & \textbf{\hspace{0pt}Jugador} & \textbf{\hspace{0pt}Cuenta} & \textbf{\hspace{0pt}Resultado} & \textbf{\hspace{0pt}Elo final} \\ \hline
\endfirsthead
\hline
\textbf{\hspace{0pt}Partida} & \textbf{\hspace{0pt}Tipo} & \textbf{\hspace{0pt}Modo} & \textbf{\hspace{0pt}Estado} & \textbf{\hspace{0pt}Jugador} & \textbf{\hspace{0pt}Cuenta} & \textbf{\hspace{0pt}Resultado} & \textbf{\hspace{0pt}Elo final} \\ \hline
\endhead
1 & CLASIFICATORIA & CLASICO & FINALIZADA & luis\_\allowbreak{}quo & REGISTRADA & GANADA & 1016 \\ \hline
1 & CLASIFICATORIA & CLASICO & FINALIZADA & marta\_\allowbreak{}muros & REGISTRADA & PERDIDA & 984 \\ \hline
2 & CLASIFICATORIA & CLASICO & FINALIZADA & marta\_\allowbreak{}muros & REGISTRADA & PERDIDA & 969 \\ \hline
2 & CLASIFICATORIA & CLASICO & FINALIZADA & luis\_\allowbreak{}quo & REGISTRADA & GANADA & 1031 \\ \hline
3 & CLASIFICATORIA & CLASICO & FINALIZADA & luis\_\allowbreak{}quo & REGISTRADA & PERDIDA & 1012 \\ \hline
3 & CLASIFICATORIA & CLASICO & FINALIZADA & marta\_\allowbreak{}muros & REGISTRADA & GANADA & 988 \\ \hline
4 & CLASIFICATORIA & CLASICO & FINALIZADA & marta\_\allowbreak{}muros & REGISTRADA & PERDIDA & 973 \\ \hline
4 & CLASIFICATORIA & CLASICO & FINALIZADA & luis\_\allowbreak{}quo & REGISTRADA & GANADA & 1027 \\ \hline
5 & CLASIFICATORIA & CLASICO & FINALIZADA & luis\_\allowbreak{}quo & REGISTRADA & GANADA & 1041 \\ \hline
5 & CLASIFICATORIA & CLASICO & FINALIZADA & marta\_\allowbreak{}muros & REGISTRADA & PERDIDA & 959 \\ \hline
6 & CLASIFICATORIA & RAPIDA & FINALIZADA & pedro\_\allowbreak{}peon & REGISTRADA & TABLAS & 1000 \\ \hline
6 & CLASIFICATORIA & RAPIDA & FINALIZADA & sofia\_\allowbreak{}salto & REGISTRADA & TABLAS & 1000 \\ \hline
7 & CLASIFICATORIA & CUATRO\_\allowbreak{}JUGADORES & FINALIZADA & luis\_\allowbreak{}quo & REGISTRADA & GANADA & 1024 \\ \hline
7 & CLASIFICATORIA & CUATRO\_\allowbreak{}JUGADORES & FINALIZADA & marta\_\allowbreak{}muros & REGISTRADA & PERDIDA & 1008 \\ \hline
7 & CLASIFICATORIA & CUATRO\_\allowbreak{}JUGADORES & FINALIZADA & pedro\_\allowbreak{}peon & REGISTRADA & PERDIDA & 976 \\ \hline
7 & CLASIFICATORIA & CUATRO\_\allowbreak{}JUGADORES & FINALIZADA & sofia\_\allowbreak{}salto & REGISTRADA & PERDIDA & 992 \\ \hline
8 & IA & CLASICO & FINALIZADA & pedro\_\allowbreak{}peon & REGISTRADA & PERDIDA & \textit{NULL} \\ \hline
9 & CLASIFICATORIA & CLASICO & FINALIZADA & luis\_\allowbreak{}quo & REGISTRADA & GANADA & 1055 \\ \hline
9 & CLASIFICATORIA & CLASICO & FINALIZADA & rayo\_\allowbreak{}veloz & REGISTRADA & PERDIDA & 986 \\ \hline
10 & PRIVADA & RAPIDA & FINALIZADA & sofia\_\allowbreak{}salto & REGISTRADA & GANADA & \textit{NULL} \\ \hline
10 & PRIVADA & RAPIDA & FINALIZADA & Invitado\_\allowbreak{}4821 & INVITADO & PERDIDA & \textit{NULL} \\ \hline
11 & CLASIFICATORIA & RAPIDA & FINALIZADA & marta\_\allowbreak{}muros & REGISTRADA & GANADA & 1016 \\ \hline
11 & CLASIFICATORIA & RAPIDA & FINALIZADA & pedro\_\allowbreak{}peon & REGISTRADA & PERDIDA & 984 \\ \hline
12 & CLASIFICATORIA & CLASICO & EN\_\allowbreak{}CURSO & luis\_\allowbreak{}quo & REGISTRADA & \textit{NULL} & \textit{NULL} \\ \hline
12 & CLASIFICATORIA & CLASICO & EN\_\allowbreak{}CURSO & sofia\_\allowbreak{}salto & REGISTRADA & \textit{NULL} & \textit{NULL} \\ \hline
\end{longtable}
\end{center}


\captura{04-relaciones}{Captura 3. La misma consulta en la consola: veinticinco participaciones de doce partidas.}

\subsubsection*{Del jugador a su aspecto}
\addcontentsline{toc}{subsubsection}{Del jugador a su aspecto}

\textit{Equipamiento} tiene dos llaves foráneas compuestas: (\at{id\_usuario}, \at{id\_objeto}) hacia \textit{UsuarioObjeto}, que obliga a equipar solo lo que se posee (\mbox{CU-14} \mbox{RN-02}), e (\at{id\_objeto}, \at{id\_ranura}) hacia \textit{ObjetoCosmetico}, que obliga a que el objeto sea de esa ranura (\mbox{D-03}). La consulta recorre la cadena completa para \at{luis\_quo}, con una fila por ranura:

\consulta{relequipamiento}
% Archivo generado por bd/exportar_datos.py. No editar a mano.
\begin{center}\footnotesize\setlength{\tabcolsep}{3pt}
\begin{longtable}{|L{0.118\linewidth}|L{0.103\linewidth}|L{0.265\linewidth}|L{0.074\linewidth}|L{0.235\linewidth}|L{0.103\linewidth}|}
\hline
\textbf{Usuario} & \textbf{Ranura} & \textbf{Objeto} & \textbf{Rareza} & \textbf{Obtenido} & \textbf{Origen} \\ \hline
\endfirsthead
\hline
\textbf{Usuario} & \textbf{Ranura} & \textbf{Objeto} & \textbf{Rareza} & \textbf{Obtenido} & \textbf{Origen} \\ \hline
\endhead
luis\_\allowbreak{}quo & PEON & PEON\_\allowbreak{}CRISTAL & RARO & 2026-07-27 18:00 & COMPRA \\ \hline
luis\_\allowbreak{}quo & MURO & MURO\_\allowbreak{}MADERA & COMUN & 2026-06-10 16:00 & INICIAL \\ \hline
luis\_\allowbreak{}quo & TABLERO & TABLERO\_\allowbreak{}CLASICO & COMUN & 2026-06-10 16:00 & INICIAL \\ \hline
luis\_\allowbreak{}quo & TITULO & TITULO\_\allowbreak{}CONSTRUCTOR & RARO & 2026-07-20 19:10 & NIVEL \\ \hline
luis\_\allowbreak{}quo & MARCO & MARCO\_\allowbreak{}SENCILLO & COMUN & 2026-06-10 16:00 & INICIAL \\ \hline
luis\_\allowbreak{}quo & EMOTES & EMOTE\_\allowbreak{}RISA & RARO & 2026-07-21 18:00 & COMPRA \\ \hline
\end{longtable}
\end{center}


Los seis objetos llegan por tres caminos distintos: los iniciales del alta (\mbox{CU-02} \mbox{RN-08}), el título que ganó al subir de nivel y los dos que compró.

\subsubsection*{Del reporte a la apelación}
\addcontentsline{toc}{subsubsection}{Del reporte a la apelación}

Un \textit{Reporte} tiene dos relaciones con \textit{Usuario}, una por papel, y una con \textit{MotivoReporte}; una \textit{Sancion} puede venir de un reporte, y una \textit{Apelacion} pertenece a una sanción. Las uniones externas conservan los reportes sin sanción y las sanciones sin apelación:

\consulta{relmoderacion}
% Archivo generado por bd/exportar_datos.py. No editar a mano.
\begin{center}\scriptsize\setlength{\tabcolsep}{3pt}
\begin{longtable}{|L{0.045\linewidth}|L{0.151\linewidth}|L{0.087\linewidth}|L{0.103\linewidth}|L{0.159\linewidth}|L{0.045\linewidth}|L{0.135\linewidth}|L{0.057\linewidth}|L{0.072\linewidth}|}
\hline
\textbf{Reporte} & \textbf{Motivo} & \textbf{Denunciante} & \textbf{Reportado} & \textbf{Estado del reporte} & \textbf{Sanción} & \textbf{Tipo} & \textbf{Apelación} & \textbf{Estado de la apelación} \\ \hline
\endfirsthead
\hline
\textbf{Reporte} & \textbf{Motivo} & \textbf{Denunciante} & \textbf{Reportado} & \textbf{Estado del reporte} & \textbf{Sanción} & \textbf{Tipo} & \textbf{Apelación} & \textbf{Estado de la apelación} \\ \hline
\endhead
1 & LENGUAJE\_\allowbreak{}OFENSIVO & luis\_\allowbreak{}quo & rayo\_\allowbreak{}veloz & RESUELTO\_\allowbreak{}CON\_\allowbreak{}SANCION & 1 & CUENTA TEMPORAL & \textit{NULL} & \textit{NULL} \\ \hline
1 & LENGUAJE\_\allowbreak{}OFENSIVO & luis\_\allowbreak{}quo & rayo\_\allowbreak{}veloz & RESUELTO\_\allowbreak{}CON\_\allowbreak{}SANCION & 2 & CUENTA TEMPORAL & 2 & PENDIENTE \\ \hline
2 & ACOSO & marta\_\allowbreak{}muros & rayo\_\allowbreak{}veloz & PENDIENTE & \textit{NULL} & \textit{NULL} & \textit{NULL} & \textit{NULL} \\ \hline
3 & JUEGO\_\allowbreak{}ANTIDEPORTIVO & sofia\_\allowbreak{}salto & luis\_\allowbreak{}quo & EN\_\allowbreak{}REVISION & \textit{NULL} & \textit{NULL} & \textit{NULL} & \textit{NULL} \\ \hline
4 & NOMBRE\_\allowbreak{}INAPROPIADO & luis\_\allowbreak{}quo & 4dm1n\_\allowbreak{}oficial & RESUELTO\_\allowbreak{}CON\_\allowbreak{}SANCION & 5 & CUENTA PERMANENTE & \textit{NULL} & \textit{NULL} \\ \hline
5 & TRAMPAS & marta\_\allowbreak{}muros & luis\_\allowbreak{}quo & RESUELTO\_\allowbreak{}SIN\_\allowbreak{}SANCION & \textit{NULL} & \textit{NULL} & \textit{NULL} & \textit{NULL} \\ \hline
\end{longtable}
\end{center}


El reporte 1 produjo dos sanciones: la segunda sustituye a la primera, y es la que rayo apeló. El 4 terminó en el baneo de \at{4dm1n\_oficial}; el 5 se resolvió sin sanción, y el 2 y el 3 siguen abiertos.

\subsubsection*{De la caja a su contenido}
\addcontentsline{toc}{subsubsection}{De la caja a su contenido}

\textit{Caja} pertenece a una cuenta y a un \textit{TipoCaja}; \textit{CajaContenido} es su entidad débil, con los tres objetos que salieron, y cada objeto repetido genera un \textit{MovimientoMoneda} de conversión que apunta a la caja y al objeto:

\consulta{relcajas}
% Archivo generado por bd/exportar_datos.py. No editar a mano.
\begin{center}\scriptsize\setlength{\tabcolsep}{3pt}
\begin{longtable}{|L{0.041\linewidth}|L{0.140\linewidth}|L{0.140\linewidth}|L{0.115\linewidth}|L{0.061\linewidth}|L{0.166\linewidth}|L{0.102\linewidth}|L{0.102\linewidth}|}
\hline
\textbf{Caja} & \textbf{Usuario} & \textbf{Tipo} & \textbf{Estado} & \textbf{Número} & \textbf{Objeto} & \textbf{Monedas por conversión} & \textbf{Movimiento} \\ \hline
\endfirsthead
\hline
\textbf{Caja} & \textbf{Usuario} & \textbf{Tipo} & \textbf{Estado} & \textbf{Número} & \textbf{Objeto} & \textbf{Monedas por conversión} & \textbf{Movimiento} \\ \hline
\endhead
1 & marta\_\allowbreak{}muros & CAJA\_\allowbreak{}MADERA & CADUCADA & \textit{NULL} & \textit{NULL} & \textit{NULL} & \textit{NULL} \\ \hline
2 & pedro\_\allowbreak{}peon & CAJA\_\allowbreak{}MADERA & SIN\_\allowbreak{}ABRIR & \textit{NULL} & \textit{NULL} & \textit{NULL} & \textit{NULL} \\ \hline
3 & luis\_\allowbreak{}quo & CAJA\_\allowbreak{}MADERA & ABIERTA & 1 & MURO\_\allowbreak{}LADRILLO & \textit{NULL} & \textit{NULL} \\ \hline
3 & luis\_\allowbreak{}quo & CAJA\_\allowbreak{}MADERA & ABIERTA & 2 & EMOTE\_\allowbreak{}RISA & 40 & 28 \\ \hline
3 & luis\_\allowbreak{}quo & CAJA\_\allowbreak{}MADERA & ABIERTA & 3 & MURO\_\allowbreak{}HIELO & \textit{NULL} & \textit{NULL} \\ \hline
\end{longtable}
\end{center}


Solo la caja abierta tiene contenido. El \at{EMOTE\_RISA} que salió en ella ya lo tenía luis, así que se convirtió en 40 monedas (\mbox{CU-39} \mbox{RN-03}) con el movimiento 28.

\subsubsection*{Una sanción que sustituye a otra}
\addcontentsline{toc}{subsubsection}{Una sanción que sustituye a otra}

La única relación de una tabla consigo misma: \at{id\_sancion\_sustituta} de \textit{Sancion} apunta a otra fila de \textit{Sancion} (\mbox{CU-44} \mbox{RN-09}). La sanción sustituida queda retirada, con su fecha de retiro, y no se borra:

\consulta{relsustitucion}
% Archivo generado por bd/exportar_datos.py. No editar a mano.
\begin{center}\footnotesize\setlength{\tabcolsep}{3pt}
\begin{longtable}{|L{0.055\linewidth}|L{0.139\linewidth}|L{0.139\linewidth}|L{0.062\linewidth}|L{0.139\linewidth}|L{0.364\linewidth}|}
\hline
\textbf{Sanción retirada} & \textbf{Fin previsto} & \textbf{Retiro} & \textbf{Sustituta} & \textbf{Fin de la sustituta} & \textbf{Motivo de la sustituta} \\ \hline
\endfirsthead
\hline
\textbf{Sanción retirada} & \textbf{Fin previsto} & \textbf{Retiro} & \textbf{Sustituta} & \textbf{Fin de la sustituta} & \textbf{Motivo de la sustituta} \\ \hline
\endhead
1 & 2026-08-28 10:00 & 2026-08-27 16:00 & 2 & 2026-09-03 16:00 & Insultos reiterados; se amplía la suspensión. \\ \hline
\end{longtable}
\end{center}


%=====================================================================
% Datos · Estados y categorías
%   La tabla de cobertura la genera bd/exportar_datos.py a partir de las
%   restricciones CHECK de la base.
%=====================================================================
\subsection*{Estados y categorías}
\addcontentsline{toc}{subsection}{Estados y categorías}

Casi todas las entidades del modelo tienen un estado, un tipo o una categoría que decide qué columnas aplican y qué puede pasarle a la fila. La captura 4 cuenta las filas de cada valor en las dieciocho columnas más importantes: las cuentas están en los cinco estados y tienen los tres roles, las partidas son de los tres tipos y terminan de las cinco formas, y los reportes, las sanciones, las cajas y las invitaciones pasan por todos sus estados.

\captura{04-estados}{Captura 4. Filas por valor en las columnas de estado y de categoría.}

La tabla siguiente hace la misma cuenta para todas las columnas de dominio cerrado: las \dnNumDominios{} cuyos valores fija una restricción \at{CHECK}. \at{bd/exportar\_datos.py} lee esos valores del catálogo del sistema, así que la tabla no puede omitir ninguno. \dnNumDominiosCompletos{} columnas tienen registros en todos sus valores, y en total \dnNumValoresConFilas{} de los \dnNumValores{} valores tienen al menos uno.

% Archivo generado por bd/exportar_datos.py. No editar a mano.
\begin{center}\footnotesize\setlength{\tabcolsep}{2pt}
\begin{longtable}{|L{143.2pt}|L{159.3pt}|L{155.2pt}|}
\hline
\textbf{\hspace{0pt}Columna} & \textbf{\hspace{0pt}Valores con filas (filas)} & \textbf{\hspace{0pt}Valores sin filas} \\ \hline
\endfirsthead
\hline
\textbf{\hspace{0pt}Columna} & \textbf{\hspace{0pt}Valores con filas (filas)} & \textbf{\hspace{0pt}Valores sin filas} \\ \hline
\endhead
ObjetoCosmetico.rareza & LEGENDARIO (2), EPICO (5), RARO (7), COMUN (9) & — \\ \hline
PalabraProhibida.ambito & AMBOS (3), CHAT (2), NICKNAME (3) & — \\ \hline
PalabraProhibida.idioma & en (3), es-MX (4), NULL (1) & — \\ \hline
Usuario.estado\_\allowbreak{}cuenta & ELIMINADA (1), BANEADA (1), SUSPENDIDA (1), ACTIVA (7), PENDIENTE (1) & — \\ \hline
Usuario.idioma\_\allowbreak{}preferido & en (1), es-MX (10) & — \\ \hline
Usuario.rol & ADMINISTRADOR (1), MODERADOR (1), JUGADOR (9) & — \\ \hline
Usuario.tipo\_\allowbreak{}cuenta & REGISTRADA (10), INVITADO (1) & — \\ \hline
Sesion.motivo\_\allowbreak{}cierre & SANCION (2), BAJA\_\allowbreak{}CUENTA (1), CAMBIO\_\allowbreak{}CREDENCIALES (1), CIERRE\_\allowbreak{}REMOTO (1), CIERRE\_\allowbreak{}VOLUNTARIO (2), NULL (8) & — \\ \hline
CodigoSegundoFactor.estado & INVALIDADO (1), CONSUMIDO (2), PENDIENTE (1) & — \\ \hline
TokenVerificacionCorreo.estado & INVALIDADO (1), USADO (9), PENDIENTE (2) & — \\ \hline
TokenVerificacionCorreo.proposito & ALTA (11), CAMBIO\_\allowbreak{}CORREO (1) & — \\ \hline
TokenRecuperacion.estado & INVALIDADO (1), USADO (1), PENDIENTE (1) & — \\ \hline
AceptacionTerminos.idioma & en (1), es-MX (10) & — \\ \hline
Avatar.formato & PNG (1), JPG (2) & — \\ \hline
UsuarioObjeto.origen & NIVEL (1), CAJA (2), COMPRA (4), INICIAL (88) & — \\ \hline
Partida.estado & FINALIZADA (11), EN\_\allowbreak{}CURSO (1) & PENDIENTE\_\allowbreak{}DE\_\allowbreak{}CIERRE \\ \hline
Partida.forma\_\allowbreak{}termino & ABANDONO (1), TABLAS (1), RENDICION (2), TIEMPO\_\allowbreak{}AGOTADO (1), META (6), NULL (1) & — \\ \hline
Partida.tipo & IA (1), PRIVADA (1), CLASIFICATORIA (10) & — \\ \hline
Participacion.forma\_\allowbreak{}termino & ABANDONO (3), NULL (22) & RENDICION, TIEMPO\_\allowbreak{}AGOTADO, META \\ \hline
Participacion.resultado & TABLAS (2), PERDIDA (12), GANADA (9), NULL (2) & — \\ \hline
Jugada.orientacion & VERTICAL (6), HORIZONTAL (17), NULL (111) & — \\ \hline
Jugada.tipo & MOVIMIENTO (111), MURO (23) & — \\ \hline
OfertaTablas.estado & CADUCADA (1), RECHAZADA (1), ACEPTADA (1), PENDIENTE (1) & — \\ \hline
Sala.estado & CERRADA (1), ABIERTA (1) & EN\_\allowbreak{}PARTIDA \\ \hline
Invitacion.estado & EXPIRADA (1), RECHAZADA (1), ACEPTADA (1), PENDIENTE (1) & — \\ \hline
Mensaje.canal & ESPECTADORES (1), PARTIDA (4), SALA (3), AMIGOS (1), GLOBAL (2) & — \\ \hline
Caja.estado & CADUCADA (1), ABIERTA (1), SIN\_\allowbreak{}ABRIR (1) & — \\ \hline
Caja.origen & PARTIDA (2), COMPRA (1) & — \\ \hline
MovimientoMoneda.tipo & CONVERSION (1), COMPRA\_\allowbreak{}CAJA (1), COMPRA (4), NIVEL (4), PARTIDA (20) & DEVOLUCION \\ \hline
Reporte.estado & PENDIENTE (1), EN\_\allowbreak{}REVISION (1), RESUELTO\_\allowbreak{}CON\_\allowbreak{}SANCION (2), RESUELTO\_\allowbreak{}SIN\_\allowbreak{}SANCION (1) & — \\ \hline
Sancion.ambito & CHAT (2), CUENTA (3) & — \\ \hline
Sancion.tipo & TEMPORAL (3), PERMANENTE (2) & — \\ \hline
Apelacion.estado & PENDIENTE (1), RECHAZADA (1), ACEPTADA (1) & EN\_\allowbreak{}REVISION \\ \hline
BitacoraModeracion.accion & INTENTO\_\allowbreak{}SIN\_\allowbreak{}PERMISO (1), CONSULTA\_\allowbreak{}BITACORA (1), ROL\_\allowbreak{}OTORGADO (1), APELACION\_\allowbreak{}RESUELTA (2), SANCION\_\allowbreak{}APLICADA (5), REPORTE\_\allowbreak{}RESUELTO (3), REPORTE\_\allowbreak{}LIBERADO (1), REPORTE\_\allowbreak{}TOMADO (5) & ROL\_\allowbreak{}RETIRADO \\ \hline
BitacoraAcceso.resultado & CUENTA\_\allowbreak{}ELIMINADA (1), SEGUNDO\_\allowbreak{}FACTOR\_\allowbreak{}ACTIVADO (1), CAMBIO\_\allowbreak{}CORREO\_\allowbreak{}SOLICITADO (1), CAMBIO\_\allowbreak{}CONTRASENA (1), CIERRE\_\allowbreak{}REMOTO (1), CIERRE\_\allowbreak{}VOLUNTARIO (2), RECUPERACION\_\allowbreak{}EXITOSA (1), RECUPERACION\_\allowbreak{}CORREO\_\allowbreak{}INEXISTENTE (1), RECUPERACION\_\allowbreak{}SOLICITADA (3), CUENTA\_\allowbreak{}VERIFICADA (9), ALTA\_\allowbreak{}INVITADO (1), REGISTRO\_\allowbreak{}RECHAZADO (1), REGISTRO\_\allowbreak{}EXITOSO (10), SEGUNDO\_\allowbreak{}FACTOR\_\allowbreak{}FALLIDO (1), CUENTA\_\allowbreak{}SANCIONADA (2), CUENTA\_\allowbreak{}PENDIENTE (1), CONTRASENA\_\allowbreak{}INCORRECTA (3), BLOQUEO\_\allowbreak{}VIGENTE (1), USUARIO\_\allowbreak{}INEXISTENTE (1), EXITO (12) & SEGUNDO\_\allowbreak{}FACTOR\_\allowbreak{}DESACTIVADO, CAMBIO\_\allowbreak{}CORREO\_\allowbreak{}CONFIRMADO, VINCULACION\_\allowbreak{}EXITOSA, REGISTRO\_\allowbreak{}SIN\_\allowbreak{}CORREO, SUSPENSION\_\allowbreak{}VENCIDA, ABANDONADO \\ \hline
\end{longtable}
\end{center}


Los valores sin registros son estados de excepción o de paso, o eventos que no ocurren en la fotografía; incluirlos obligaría a romper la coherencia del escenario:

\begin{reglas}
  \item \textbf{\at{PENDIENTE\_DE\_CIERRE} de \textit{Partida}.} Una partida queda así solo si se agotan los reintentos al cerrarla (\mbox{CU-25} \mbox{EX-01}), y el servidor la cierra en cuanto puede.
  \item \textbf{\at{EN\_PARTIDA} de \textit{Sala}.} Una sala lo está solo mientras dura su partida privada (\mbox{CU-18}). La de la sala 1 ya terminó, y la única partida en curso es clasificatoria.
  \item \textbf{\at{META}, \at{TIEMPO\_AGOTADO} y \at{RENDICION} en la salida de una \textit{Participacion}.} Esa columna solo se escribe cuando un jugador sale antes que los demás de una partida de cuatro, y en la de prueba los tres salen por abandono. Las cinco formas sí aparecen como forma de término de las partidas.
  \item \textbf{\at{EN\_REVISION} de \textit{Apelacion}.} Una apelación tomada vuelve a la cola a los treinta minutos (\mbox{CU-43} \mbox{RN-03}), y a la hora de la fotografía la única moderadora conectada es la que aplicó las sanciones apeladas, que no puede resolverlas (\mbox{CU-45} \mbox{RN-06}).
  \item \textbf{\at{DEVOLUCION} de \textit{MovimientoMoneda}.} Solo aparece si falla la confirmación de la compra de una caja (\mbox{CU-39} \mbox{EX-02}) o para corregir un movimiento (\mbox{CU-40} \mbox{RN-01}), y en el escenario no falla ninguna compra.
  \item \textbf{Acciones de las bitácoras.} Nadie retira un rol, desactiva el segundo factor, confirma un cambio de correo, vincula su cuenta de invitado ni abandona un inicio de sesión a medias; el alta no encuentra caído el servicio de correo, y la única suspensión vigente aún no vence.
\end{reglas}

%=====================================================================
% Datos · Variaciones
%=====================================================================
\subsection*{Variaciones}
\addcontentsline{toc}{subsection}{Variaciones}

Además de los estados, los datos incluyen registros que se salen del caso normal en la forma de su información: columnas que quedan nulas según el tipo de fila, textos que no están en español o que llevan emojis, partidas sin reloj o de cuatro jugadores, registros que no se borran aunque dejen de valer. La consulta siguiente los reúne con una etiqueta que dice qué variación representa cada uno:

\consulta{variaciones}
% Archivo generado por bd/exportar_datos.py. No editar a mano.
\begin{center}\footnotesize\setlength{\tabcolsep}{2pt}
\begin{longtable}{|L{130.6pt}|L{86.3pt}|L{94.5pt}|L{142.0pt}|}
\hline
\textbf{\hspace{0pt}Variación} & \textbf{\hspace{0pt}Tabla} & \textbf{\hspace{0pt}Registro} & \textbf{\hspace{0pt}Dato que la muestra} \\ \hline
\endfirsthead
\hline
\textbf{\hspace{0pt}Variación} & \textbf{\hspace{0pt}Tabla} & \textbf{\hspace{0pt}Registro} & \textbf{\hspace{0pt}Dato que la muestra} \\ \hline
\endhead
Invitado: sin correo ni contraseña & Usuario & Invitado\_\allowbreak{}4821 & correo NULL, contraseña NULL \\ \hline
Cuenta eliminada y anonimizada (D-17) & Usuario & anonimo\_\allowbreak{}10 & correo anonimo\_\allowbreak{}10, baja 2026-08-02 13:00 \\ \hline
Cuenta baneada: sanción permanente & Sancion & 4dm1n\_\allowbreak{}oficial & CUENTA PERMANENTE, fin NULL \\ \hline
Interfaz y términos en inglés (D-21) & AceptacionTerminos & sofia\_\allowbreak{}salto & versión 2026.1 en en \\ \hline
Chat con ñ, acentos, emoji e inglés & Mensaje & mensaje 4 & Welcome! We can start whenever you want \\ \hline
Chat con ñ, acentos, emoji e inglés & Mensaje & mensaje 7 & Igualmente, compañera 👍 \\ \hline
Partida contra la IA, sin reloj & Partida & partida 8 & reloj NULL, deshacer usados 1 \\ \hline
Jugadas deshechas que se conservan (D-12) & Jugada & partida 8 & jugadas 5, 6 con deshecha = 1 \\ \hline
Cuatro jugadores: posiciones y salidas & Participacion & luis\_\allowbreak{}quo & posición 1, salida ninguna \\ \hline
Cuatro jugadores: posiciones y salidas & Participacion & marta\_\allowbreak{}muros & posición 2, salida ABANDONO \\ \hline
Cuatro jugadores: posiciones y salidas & Participacion & pedro\_\allowbreak{}peon & posición 4, salida ABANDONO \\ \hline
Cuatro jugadores: posiciones y salidas & Participacion & sofia\_\allowbreak{}salto & posición 3, salida ABANDONO \\ \hline
Derrota por tiempo: reloj en cero & Participacion & pedro\_\allowbreak{}peon & partida 11, reloj 0 ms \\ \hline
Objeto retirado del catálogo & ObjetoCosmetico & PEON\_\allowbreak{}TEMPORADA\_\allowbreak{}1 & activo 0, a la venta 0 \\ \hline
Término prohibido en todos los idiomas & PalabraProhibida & admin & idioma NULL \\ \hline
Sesión que expiró sin cierre & Sesion & sesión 12 & expiró 2026-08-31 16:00, fin NULL \\ \hline
Sesión que expiró sin cierre & Sesion & sesión 14 & expiró 2026-09-01 16:20, fin NULL \\ \hline
\end{longtable}
\end{center}


\captura{05-variaciones}{Captura 5. Las mismas variaciones en la consola. La consola dibuja el emoji del mensaje 7 como un recuadro porque su fuente no lo tiene.}

Lo que cada grupo permite probar:

\begin{reglas}
  \item \textbf{Nulos que dependen del tipo de fila.} El invitado no tiene correo ni contraseña, y la cuenta eliminada conserva su fila con los datos personales anonimizados (\mbox{D-07}, \mbox{D-17}). Las restricciones \at{CHECK} de subtipo los exigen o los prohíben según el tipo de cuenta.
  \item \textbf{Sanciones permanentes.} La sanción permanente no tiene fecha de fin; la temporal debe tenerla (\mbox{CU-44} \mbox{RN-04}), como comprueba la prueba 11 de los casos que la base rechaza.
  \item \textbf{Idioma y codificación.} \at{sofia\_salto} aceptó los términos en inglés y escribe en inglés; el mensaje 7 lleva una eñe y un emoji. La última comprobación del script de datos verifica que esos textos se guardaron sin pérdida (\mbox{D-21}).
  \item \textbf{Partidas fuera de lo común.} La partida contra la IA no tiene reloj y conserva las dos jugadas que se deshicieron; la de cuatro jugadores guarda la posición de cada uno y la forma en que salió; la perdida por tiempo deja el reloj del perdedor en cero.
  \item \textbf{Lo que deja de valer sin borrarse.} El objeto retirado sigue en el catálogo, sin venderse, porque hay quien lo tiene; las sesiones que expiraron sin cerrarse siguen en su tabla, y el término prohibido sin idioma se filtra en todos.
\end{reglas}

Las jugadas de la partida contra la IA muestran cómo se guarda un deshacer: las jugadas 5 y 6 quedan marcadas como deshechas y la partida sigue desde la posición anterior a ellas, así que en la 7 pedro mueve a otra casilla y en la 8 la IA repite su movimiento (\mbox{D-12}). Las jugadas de la IA no tienen usuario.

\consulta{variaia}
% Archivo generado por bd/exportar_datos.py. No editar a mano.
\begin{center}\scriptsize\setlength{\tabcolsep}{2pt}
\begin{longtable}{|L{50.0pt}|L{54.9pt}|L{84.4pt}|L{42.7pt}|L{46.8pt}|L{35.7pt}|L{85.1pt}|L{36.2pt}|}
\hline
\textbf{\hspace{0pt}Número} & \textbf{\hspace{0pt}Jugador} & \textbf{\hspace{0pt}Tipo} & \textbf{\hspace{0pt}Origen} & \textbf{\hspace{0pt}Destino} & \textbf{\hspace{0pt}Surco} & \textbf{\hspace{0pt}Orientación} & \textbf{\hspace{0pt}Deshecha} \\ \hline
\endfirsthead
\hline
\textbf{\hspace{0pt}Número} & \textbf{\hspace{0pt}Jugador} & \textbf{\hspace{0pt}Tipo} & \textbf{\hspace{0pt}Origen} & \textbf{\hspace{0pt}Destino} & \textbf{\hspace{0pt}Surco} & \textbf{\hspace{0pt}Orientación} & \textbf{\hspace{0pt}Deshecha} \\ \hline
\endhead
1 & pedro\_\allowbreak{}peon & MOVIMIENTO & e1 & e2 & \textit{NULL} & \textit{NULL} & 0 \\ \hline
2 & IA & MOVIMIENTO & e9 & e8 & \textit{NULL} & \textit{NULL} & 0 \\ \hline
3 & pedro\_\allowbreak{}peon & MOVIMIENTO & e2 & e3 & \textit{NULL} & \textit{NULL} & 0 \\ \hline
4 & IA & MURO & \textit{NULL} & \textit{NULL} & e4 & HORIZONTAL & 0 \\ \hline
5 & pedro\_\allowbreak{}peon & MOVIMIENTO & e3 & d3 & \textit{NULL} & \textit{NULL} & 1 \\ \hline
6 & IA & MOVIMIENTO & e8 & e7 & \textit{NULL} & \textit{NULL} & 1 \\ \hline
7 & pedro\_\allowbreak{}peon & MOVIMIENTO & e3 & f3 & \textit{NULL} & \textit{NULL} & 0 \\ \hline
8 & IA & MOVIMIENTO & e8 & e7 & \textit{NULL} & \textit{NULL} & 0 \\ \hline
9 & pedro\_\allowbreak{}peon & MOVIMIENTO & f3 & f4 & \textit{NULL} & \textit{NULL} & 0 \\ \hline
\end{longtable}
\end{center}


%=====================================================================
% Datos · Casos que la base rechaza
%=====================================================================
\subsection*{Casos que la base rechaza}
\addcontentsline{toc}{subsection}{Casos que la base rechaza}

Unos datos de prueba solo sirven si la base también rechaza los que no deberían existir. \at{bd/pruebas\_rechazo.sql} intenta, sobre los registros cargados, \dnNumPruebasRechazo{} operaciones que el documento prohíbe, cada una en su propia transacción, que se revierte siempre: el script no cambia ningún dato. De cada intento anota el número de error de SQL Server y la restricción que lo impidió, que sale del propio mensaje de error.

% Archivo generado por bd/exportar_datos.py. No editar a mano.
\begin{center}\footnotesize\setlength{\tabcolsep}{2pt}
\begin{longtable}{|L{29.5pt}|L{166.5pt}|L{83.0pt}|L{26.1pt}|L{143.9pt}|}
\hline
\textbf{\hspace{0pt}Núm.} & \textbf{\hspace{0pt}Operación que se intenta} & \textbf{\hspace{0pt}Regla} & \textbf{\hspace{0pt}Error} & \textbf{\hspace{0pt}Rechazada por} \\ \hline
\endfirsthead
\hline
\textbf{\hspace{0pt}Núm.} & \textbf{\hspace{0pt}Operación que se intenta} & \textbf{\hspace{0pt}Regla} & \textbf{\hspace{0pt}Error} & \textbf{\hspace{0pt}Rechazada por} \\ \hline
\endhead
1 & Estado de cuenta fuera del dominio & CRUD, dominios & 547 & CK\_\allowbreak{}Usuario\_\allowbreak{}estado\_\allowbreak{}cuenta \\ \hline
2 & Nickname igual salvo mayúsculas y acentos & CU-02 RN-01 & 2627 & UQ\_\allowbreak{}Usuario\_\allowbreak{}nickname \\ \hline
3 & Idioma no admitido & CU-12 FA-07, D-21 & 547 & CK\_\allowbreak{}Usuario\_\allowbreak{}idioma \\ \hline
4 & Saldo de monedas negativo & CU-40 RN-02 & 547 & CK\_\allowbreak{}Usuario\_\allowbreak{}rangos \\ \hline
5 & Amistad con una cuenta que no existe & Llave foránea & 547 & FK\_\allowbreak{}Amistad\_\allowbreak{}UsuarioB \\ \hline
6 & Equipar un objeto que no se posee & CU-14 RN-02 & 547 & FK\_\allowbreak{}Equipamiento\_\allowbreak{}UsuarioObjeto \\ \hline
7 & Equipar un objeto en otra ranura & D-03, CU-14 RN-01 & 547 & FK\_\allowbreak{}Equipamiento\_\allowbreak{}ObjetoCosmetico \\ \hline
8 & Segunda partida sin terminar & D-11 & 2601 & UX\_\allowbreak{}Participacion\_\allowbreak{}sin\_\allowbreak{}terminar \\ \hline
9 & Segundo código de recuperación vigente & CU-03 RN-05 & 2601 & UX\_\allowbreak{}TokenRecuperacion\_\allowbreak{}vigente \\ \hline
10 & Reportarse a sí mismo & CU-42 FA-06 & 547 & CK\_\allowbreak{}Reporte\_\allowbreak{}distintos \\ \hline
11 & Sanción temporal sin fecha de fin & CU-44 RN-04 & 547 & CK\_\allowbreak{}Sancion\_\allowbreak{}tipo \\ \hline
12 & Compra que no dice qué objeto & CU-38 RN-08 & 547 & CK\_\allowbreak{}MovimientoMoneda\_\allowbreak{}origen \\ \hline
13 & Borrar con el usuario de conexión & Permisos de conexión & 229 & permiso DELETE denegado \\ \hline
\end{longtable}
\end{center}


\captura{06-rechazos}{Captura 6. Las pruebas de rechazo en la consola.}

La base rechazó las \dnNumPruebasRechazadas{} operaciones, y cada una con la restricción que corresponde a su regla. El error 547 es el de una restricción \at{CHECK} o de una llave foránea; el 2627, el de una restricción \at{UNIQUE}; el 2601, el de un índice único filtrado, y el 229, el de un permiso denegado. Tres resultados merecen un comentario:

\begin{reglas}
  \item \textbf{Prueba 2.} \at{LUIS\_QUÓ} se considera el mismo nickname que \at{luis\_quo}, porque la columna compara sin distinguir mayúsculas ni acentos (\mbox{CU-02} \mbox{RN-01}).
  \item \textbf{Pruebas 6 y 7.} Las dos las rechazan las llaves foráneas compuestas de \textit{Equipamiento}: la primera, porque pedro no posee el objeto; la segunda, porque el objeto es un muro y la ranura es la del peón.
  \item \textbf{Prueba 13.} El usuario de conexión del servidor no puede borrar, ni siquiera un silencio, que es de las pocas filas que los casos de uso eliminan: para eso tiene los procedimientos almacenados.
\end{reglas}

%=====================================================================
% Datos · Registros por tabla
%   Cada tabla la genera bd/exportar_datos.py con la consulta del mismo
%   nombre de bd/consultas_datos.sql.
%=====================================================================
\subsection*{Registros por tabla}
\addcontentsline{toc}{subsection}{Registros por tabla}

% \registros{consulta}{título}{comentario}
\newcommand{\registros}[3]{\par\needspace{6\baselineskip}\noindent\textbf{#2.} #3\input{bd/tex/datos/#1}}

El contenido de las tablas, agrupadas por área como en el modelo de datos. Las llaves foráneas se muestran con el nickname o el código de la fila a la que apuntan, y los resúmenes criptográficos de contraseñas y tokens se omiten. Cuatro tablas se resumen porque sus filas se deducen de otras: \textit{Jugada}, por partida, con las jugadas completas de las partidas 1 y 8 en las subsecciones anteriores; \textit{UsuarioObjeto}, por cuenta y origen; \textit{EstadisticaModo}, solo con las filas de los modos en que la cuenta ha jugado partidas clasificatorias; y \textit{ParticipacionObjeto}, que no se muestra: sus filas repiten, para cada participación, lo que el jugador llevaba equipado al empezar (\mbox{CU-14} \mbox{RN-03}). Todos los registros, sin resumir, están en el listado de \at{insertar\_datos\_prueba.sql} al final de la sección de implementación.

\subsubsection*{Catálogos precargados}
\addcontentsline{toc}{subsubsection}{Catálogos precargados}

\registros{ranura}{\textit{Ranura}}{Las seis ranuras de personalización (\mbox{D-03}).}
\registros{modo}{\textit{Modo}}{Los tres modos con su tablero, sus plazas y sus muros (\mbox{CU-21} \mbox{RN-02}).}
\registros{division}{\textit{Division}}{En orden de \at{elo\_minimo}.}
\registros{nivelia}{\textit{NivelIA}}{Los cuatro niveles de la IA con su fuerza (\mbox{CU-27} \mbox{RN-02}).}
\registros{leccion}{\textit{Leccion}}{Las siete lecciones del tutorial (\mbox{CU-41} \mbox{RN-01}).}
\registros{tipocaja}{\textit{TipoCaja}}{La de temporada ya no se vende.}
\registros{motivo}{\textit{MotivoReporte}}{Los cinco motivos de reporte (\mbox{CU-42} \mbox{RN-08}).}
\registros{objeto}{\textit{ObjetoCosmetico}}{Uno predeterminado por ranura; el de temporada, retirado.}
\registros{tipocajaobjeto}{\textit{TipoCajaObjeto}}{Las probabilidades de cada tipo de caja suman 100 (\mbox{CU-39} \mbox{RN-01}).}
\registros{palabra}{\textit{PalabraProhibida}}{Sin idioma, el término se filtra en todos (\mbox{CU-28} \mbox{RN-01}).}

\subsubsection*{Cuenta y acceso}
\addcontentsline{toc}{subsubsection}{Cuenta y acceso}

\registros{usuario}{\textit{Usuario}}{Las once cuentas del escenario.}

\captura{07-usuarios}{Captura 7. Las cuentas en la consola.}

\registros{sesion}{\textit{Sesion}}{Las que tienen fin nulo siguen abiertas o expiraron sin cerrarse.}
\registros{aceptacion}{\textit{AceptacionTerminos}}{Desde el 1 de agosto rige la versión 2026.2.}
\registros{tokenverificacion}{\textit{TokenVerificacionCorreo}}{}
\registros{tokenrecuperacion}{\textit{TokenRecuperacion}}{}
\registros{codigo2fa}{\textit{CodigoSegundoFactor}}{}

\subsubsection*{Perfil y personalización}
\addcontentsline{toc}{subsubsection}{Perfil y personalización}

\registros{perfil}{\textit{HistorialNickname}, \textit{Avatar} y \textit{EnlaceRed}}{}
\registros{usuarioobjeto}{\textit{UsuarioObjeto}}{Objetos de cada cuenta según cómo los obtuvo.}
\registros{equipamiento}{\textit{Equipamiento}}{Lo que cada cuenta lleva en cada ranura.}

\subsubsection*{Partida}
\addcontentsline{toc}{subsubsection}{Partida}

\registros{partida}{\textit{Partida}}{En la partida contra la IA, el tipo va seguido del nivel de la IA.}
\registros{participacion}{\textit{Participacion}}{El reloj restante está en milisegundos.}
\registros{jugadasporpartida}{\textit{Jugada}}{Resumen por partida.}
\registros{oferta}{\textit{OfertaTablas}}{}
\registros{enlaceespectador}{\textit{EnlaceEspectador}}{}

\subsubsection*{Emparejamiento, salas y chat}
\addcontentsline{toc}{subsubsection}{Emparejamiento, salas y chat}

\registros{cola}{\textit{ColaEmparejamiento}}{}
\registros{sala}{\textit{Sala}}{}
\registros{salaparticipante}{\textit{SalaParticipante}}{}
\registros{invitacion}{\textit{Invitacion}}{}
\registros{mensaje}{\textit{Mensaje}}{}

\subsubsection*{Clasificación y economía}
\addcontentsline{toc}{subsubsection}{Clasificación y economía}

\registros{estadistica}{\textit{EstadisticaModo}}{Filas de los modos en que la cuenta ha jugado partidas clasificatorias.}
\registros{historialdivision}{\textit{HistorialDivision}}{}
\registros{caja}{\textit{Caja}}{}
\registros{cajacontenido}{\textit{CajaContenido}}{}
\registros{movimiento}{\textit{MovimientoMoneda}}{El saldo de cada cuenta es la suma de sus movimientos (\mbox{CU-40} \mbox{RN-02}).}

\subsubsection*{Relaciones entre jugadores y tutorial}
\addcontentsline{toc}{subsubsection}{Relaciones entre jugadores y tutorial}

\registros{social}{\textit{Amistad}, \textit{Solicitud}, \textit{Silencio} y \textit{Bloqueo}}{}
\registros{tutorial}{\textit{ProgresoTutorial}}{}

\subsubsection*{Moderación y bitácoras}
\addcontentsline{toc}{subsubsection}{Moderación y bitácoras}

\registros{reporte}{\textit{Reporte}}{}
\registros{sancion}{\textit{Sancion}}{}
\registros{apelacion}{\textit{Apelacion}}{}
\registros{bitacoramoderacion}{\textit{BitacoraModeracion}}{}
\registros{bitacoraacceso}{\textit{BitacoraAcceso}}{}

